\section{Matrix Representation}

\begin{frame}{Matrix Form for Efficiency}
  \begin{columns}
    \begin{column}{0.65\textwidth}
      \begin{mathbox}{Cubic Bézier Matrix Form}

        For a cubic Bézier curve:{
          \scriptsize
          \begin{align*}
            \mathbf{B}(t) =
            \begin{bmatrix} 1 & t & t^2 & t^3
            \end{bmatrix}
            \begin{bmatrix}
              1 & 0 & 0 & 0 \\
              -3 & 3 & 0 & 0 \\
              3 & -6 & 3 & 0 \\
              -1 & 3 & -3 & 1
            \end{bmatrix}
            \begin{bmatrix}
              \mathbf{P_0} \\
              \mathbf{P_1} \\
              \mathbf{P_2} \\
              \mathbf{P_3}
            \end{bmatrix}
          \end{align*}

        }
        \vspace{0.3cm}
        Or more compactly:
        \begin{align*}
          \mathbf{B}(t) = \mathbf{T}(t) \cdot \mathbf{M_B} \cdot \mathbf{P}
        \end{align*}
      \end{mathbox}
    \end{column}
    \begin{column}{0.35\textwidth}
      \begin{conceptbox}{Matrix Components}
        \scriptsize
        \textbf{$\mathbf{T}(t)$:} Parameter vector
        \begin{align*}
          \mathbf{T}(t) =
          \begin{bmatrix} 1 & t & t^2 & t^3
          \end{bmatrix}
        \end{align*}

        \textbf{$\mathbf{M_B}$:} Bézier basis matrix

        \textbf{$\mathbf{P}$:} Control point vector
        \begin{align*}
          \mathbf{P} =
          \begin{bmatrix} \mathbf{P_0} \\ \mathbf{P_1} \\ \mathbf{P_2} \\ \mathbf{P_3}
          \end{bmatrix}
        \end{align*}
      \end{conceptbox}
    \end{column}
  \end{columns}
\end{frame}

\begin{frame}[fragile]{Why Use Matrix Form?}
  \begin{columns}
    \begin{column}{0.5\textwidth}
      \begin{conceptbox}{Computational Advantages}
        \footnotesize
        \textbf{1. Efficiency:}
        \begin{itemize}
          \item Single matrix multiplication
          \item GPU-friendly
        \end{itemize}

        \textbf{2. Derivatives:}
        \begin{itemize}
          \item Easy to compute tangents
          \item First derivative: $\mathbf{B}'(t) = \mathbf{T}'(t) \cdot \mathbf{M_B} \cdot \mathbf{P}$
          \item Where $\mathbf{T}'(t) =
            \begin{bmatrix} 0 & 1 & 2t & 3t^2
            \end{bmatrix}$
        \end{itemize}

        \textbf{3. Transformations:}
        \begin{itemize}
          \item Apply transforms to control points
          \item Curve transforms automatically
        \end{itemize}
      \end{conceptbox}
    \end{column}
    \begin{column}{0.5\textwidth}
      \begin{conceptbox}{Implementation Example}
        \scriptsize
        \begin{minted}{python}
import numpy as np

def bezier_cubic(t, control_points):
    # Parameter vector
    T = np.array([1, t, t**2, t**3])

    # Bézier basis matrix
    M = np.array([
        [ 1,  0,  0,  0],
        [-3,  3,  0,  0],
        [ 3, -6,  3,  0],
        [-1,  3, -3,  1]
    ])

    # Control points (4x2 for 2D)
    P = np.array(control_points)

    return T @ M @ P
        \end{minted}
      \end{conceptbox}
    \end{column}
  \end{columns}
\end{frame}
